\documentclass[a4paper,12pt,oneside]{memoir}

% Castellano
\usepackage[spanish,es-tabla]{babel}
\selectlanguage{spanish}
\usepackage[utf8]{inputenc}
\usepackage[T1]{fontenc}
\usepackage{lmodern} % Scalable font
\usepackage{microtype}
\usepackage{placeins}
\usepackage{float}
\usepackage{pdflscape}
\usepackage{tabularx}
\usepackage{longtable}

\RequirePackage{booktabs}
\RequirePackage[table]{xcolor}
\RequirePackage{xtab}
\RequirePackage{multirow}

% Links
\PassOptionsToPackage{hyphens}{url}\usepackage[colorlinks]{hyperref}
\hypersetup{
	colorlinks = true,
	linkcolor = {blue!45!black},
	citecolor = {blue!45!black},
	urlcolor = {blue!55!black},
	pdftitle = {Aplicación web de descarga centralizada de software con gestión de bundles y búsqueda semántica mediante embeddings},
	pdfauthor = {José Gallardo Caballero}
}
\urlstyle{same}

% Ecuaciones
\usepackage{amsmath}

% Rutas de fichero / paquete
\newcommand{\ruta}[1]{{\sffamily #1}}

% Párrafos
\nonzeroparskip

% Huérfanas y viudas
\widowpenalty100000
\clubpenalty100000
\setlength{\headheight}{16pt}
\nouppercaseheads

% Imágenes

% Comando para insertar una imagen en un lugar concreto.
% Los parámetros son:
% 1 --> Ruta absoluta/relativa de la figura
% 2 --> Texto a pie de figura
% 3 --> Tamaño en tanto por uno relativo al ancho de página
\usepackage{graphicx}
\newcommand{\imagen}[3]{
	\begin{figure}[!h]
		\centering
		\includegraphics[width=#3\textwidth]{#1}
		\caption{#2}\label{fig:#1}
	\end{figure}
	\FloatBarrier
}

% Comando para insertar una imagen sin posición.
% Los parámetros son:
% 1 --> Ruta absoluta/relativa de la figura
% 2 --> Texto a pie de figura
% 3 --> Tamaño en tanto por uno relativo al ancho de página
\newcommand{\imagenflotante}[3]{
	\begin{figure}
		\centering
		\includegraphics[width=#3\textwidth]{#1}
		\caption{#2}\label{fig:#1}
	\end{figure}
}

% Diagrama en página apaisada.
\newcommand{\diagramaapaisado}[3]{
	\begin{landscape}
		\begin{figure}[p]
			\centering
			\includegraphics[width=0.96\linewidth,height=0.86\textheight,keepaspectratio]{#1}
			\caption{#2}
			\label{#3}
		\end{figure}
	\end{landscape}
	\FloatBarrier
}

% El comando \figura nos permite insertar figuras comodamente, y utilizando
% siempre el mismo formato. Los parametros son:
% 1 --> Porcentaje del ancho de página que ocupará la figura (de 0 a 1)
% 2 --> Fichero de la imagen
% 3 --> Texto a pie de imagen
% 4 --> Etiqueta (label) para referencias
% 5 --> Opciones que queramos pasarle al \includegraphics
% 6 --> Opciones de posicionamiento a pasarle a \begin{figure}
\newcommand{\figuraConPosicion}[6]{%
  \setlength{\anchoFloat}{#1\textwidth}%
  \addtolength{\anchoFloat}{-4\fboxsep}%
  \setlength{\anchoFigura}{\anchoFloat}%
  \begin{figure}[#6]
    \begin{center}%
      \Ovalbox{%
        \begin{minipage}{\anchoFloat}%
          \begin{center}%
            \includegraphics[width=\anchoFigura,#5]{#2}%
            \caption{#3}%
            \label{#4}%
          \end{center}%
        \end{minipage}
      }%
    \end{center}%
  \end{figure}%
}

%
% Comando para incluir imágenes en formato apaisado (sin marco).
\newcommand{\figuraApaisadaSinMarco}[5]{%
  \begin{figure}%
    \begin{center}%
    \includegraphics[angle=90,height=#1\textheight,#5]{#2}%
    \caption{#3}%
    \label{#4}%
    \end{center}%
  \end{figure}%
}
% Para las tablas
\newcommand{\otoprule}{\midrule [\heavyrulewidth]}
%
% Nuevo comando para tablas pequeñas (menos de una página).
\newcommand{\tablaSmall}[5]{%
 \begin{table}
  \begin{center}
   \rowcolors {2}{gray!35}{}
   \begin{tabular}{#2}
    \toprule
    #4
    \otoprule
    #5
    \bottomrule
   \end{tabular}
   \caption{#1}
   \label{tabla:#3}
  \end{center}
 \end{table}
}

%
% Nuevo comando para tablas pequeñas (menos de una página).
\newcommand{\tablaSmallSinColores}[5]{%
 \begin{table}[H]
  \begin{center}
   \begin{tabular}{#2}
    \toprule
    #4
    \otoprule
    #5
    \bottomrule
   \end{tabular}
   \caption{#1}
   \label{tabla:#3}
  \end{center}
 \end{table}
}

\newcommand{\tablaApaisadaSmall}[5]{%
\begin{landscape}
  \begin{table}
   \begin{center}
    \rowcolors {2}{gray!35}{}
    \begin{tabular}{#2}
     \toprule
     #4
     \otoprule
     #5
     \bottomrule
    \end{tabular}
    \caption{#1}
    \label{tabla:#3}
   \end{center}
  \end{table}
\end{landscape}
}

%
% Nuevo comando para tablas grandes con cabecera y filas alternas coloreadas en gris.
\newcommand{\tabla}[6]{%
  \begin{center}
    \tablefirsthead{
      \toprule
      #5
      \otoprule
    }
    \tablehead{
      \multicolumn{#3}{l}{\small\sl continúa desde la página anterior}\\
      \toprule
      #5
      \otoprule
    }
    \tabletail{
      \hline
      \multicolumn{#3}{r}{\small\sl continúa en la página siguiente}\\
    }
    \tablelasttail{
      \hline
    }
    \bottomcaption{#1}
    \rowcolors {2}{gray!35}{}
    \begin{xtabular}{#2}
      #6
      \bottomrule
    \end{xtabular}
    \label{tabla:#4}
  \end{center}
}

%
% Nuevo comando para tablas grandes con cabecera.
\newcommand{\tablaSinColores}[6]{%
  \begin{center}
    \tablefirsthead{
      \toprule
      #5
      \otoprule
    }
    \tablehead{
      \multicolumn{#3}{l}{\small\sl continúa desde la página anterior}\\
      \toprule
      #5
      \otoprule
    }
    \tabletail{
      \hline
      \multicolumn{#3}{r}{\small\sl continúa en la página siguiente}\\
    }
    \tablelasttail{
      \hline
    }
    \bottomcaption{#1}
    \begin{xtabular}{#2}
      #6
      \bottomrule
    \end{xtabular}
    \label{tabla:#4}
  \end{center}
}

%
% Nuevo comando para tablas grandes sin cabecera.
\newcommand{\tablaSinCabecera}[5]{%
  \begin{center}
    \tablefirsthead{
      \toprule
    }
    \tablehead{
      \multicolumn{#3}{l}{\small\sl continúa desde la página anterior}\\
      \hline
    }
    \tabletail{
      \hline
      \multicolumn{#3}{r}{\small\sl continúa en la página siguiente}\\
    }
    \tablelasttail{
      \hline
    }
    \bottomcaption{#1}
  \begin{xtabular}{#2}
    #5
   \bottomrule
  \end{xtabular}
  \label{tabla:#4}
  \end{center}
}



\definecolor{cgoLight}{HTML}{EEEEEE}
\definecolor{cgoExtralight}{HTML}{FFFFFF}

%
% Nuevo comando para tablas grandes sin cabecera.
\newcommand{\tablaSinCabeceraConBandas}[5]{%
  \begin{center}
    \tablefirsthead{
      \toprule
    }
    \tablehead{
      \multicolumn{#3}{l}{\small\sl continúa desde la página anterior}\\
      \hline
    }
    \tabletail{
      \hline
      \multicolumn{#3}{r}{\small\sl continúa en la página siguiente}\\
    }
    \tablelasttail{
      \hline
    }
    \bottomcaption{#1}
    \rowcolors[]{1}{cgoExtralight}{cgoLight}

  \begin{xtabular}{#2}
    #5
   \bottomrule
  \end{xtabular}
  \label{tabla:#4}
  \end{center}
}



\graphicspath{ {./img/} }

% Capítulos
\chapterstyle{bianchi}
\newcommand{\capitulo}[2]{
	\setcounter{chapter}{#1}
	\setcounter{section}{0}
	\setcounter{figure}{0}
	\setcounter{table}{0}
	\chapter*{\thechapter.\enskip #2}
	\addcontentsline{toc}{chapter}{\thechapter.\enskip #2}
	\markboth{#2}{#2}
}

% Apéndices
\renewcommand{\appendixname}{Apéndice}
\renewcommand*\cftappendixname{\appendixname}

\newcommand{\apendice}[1]{
	%\renewcommand{\thechapter}{A}
	\chapter{#1}
}

\renewcommand*\cftappendixname{\appendixname\ }

% Formato de portada
\makeatletter
\usepackage{xcolor}
\newcommand{\tutor}[1]{\def\@tutor{#1}}
\newcommand{\course}[1]{\def\@course{#1}}
\definecolor{cpardoBox}{HTML}{E6E6FF}
\def\maketitle{
  \null
  \thispagestyle{empty}
  % Cabecera ----------------
\noindent\includegraphics[width=\textwidth]{cabecera}\vspace{1cm}%
  \vfill
  % Título proyecto y escudo informática ----------------
  \colorbox{cpardoBox}{%
    \begin{minipage}{.8\textwidth}
      \vspace{.5cm}\Large
      \begin{center}
      \textbf{TFG del Grado en Ingeniería Informática}\vspace{.6cm}\\
      \textbf{\LARGE\@title{}}
      \end{center}
      \vspace{.2cm}
    \end{minipage}

  }%
  \hfill\begin{minipage}{.20\textwidth}
    \includegraphics[width=\textwidth]{escudoInfor}
  \end{minipage}
  \vfill
  % Datos de alumno, curso y tutores ------------------
  \begin{center}%
  {%
    \noindent\LARGE
    Presentado por \@author{}\\ 
    en Universidad de Burgos --- \@date{}\\
    Tutores: \@tutor{}\\
  }%
  \end{center}%
  \null
  \cleardoublepage
  }
\makeatother

\newcommand{\nombre}{José Gallardo Caballero}

% Datos de portada
\title{Aplicación web de descarga centralizada de software con gestión de bundles y búsqueda semántica mediante embeddings}
\author{\nombre}
\tutor{José Manuel Aroca Fernández, Rodrigo Pascual García}
\date{\today}

\begin{document}

\maketitle


\newpage\null\thispagestyle{empty}\newpage


%%%%%%%%%%%%%%%%%%%%%%%%%%%%%%%%%%%%%%%%%%%%%%%%%%%%%%%%%%%%%%%%%%%%%%%%%%%%%%%%%%%%%%%%
\thispagestyle{empty}


\noindent\includegraphics[width=\textwidth]{cabecera}\vspace{1cm}

\noindent D. José Manuel Aroca Fernández, profesor del Departamento de Ingeniería Informática de la Universidad de Burgos,

\noindent Expone:

\noindent Que el alumno D. \nombre, con DNI 71478806G, ha realizado el Trabajo final de Grado en Ingeniería Informática titulado \emph{Aplicación web de descarga centralizada de software con gestión de bundles y búsqueda semántica mediante embeddings}.

\noindent Y que dicho trabajo ha sido realizado por el alumno bajo la dirección del que suscribe, en virtud de lo cual se autoriza su presentación y defensa.

\begin{center} %\large
En Burgos, {\large \today}
\end{center}

\vfill\vfill\vfill

% Author and supervisor
\begin{minipage}{0.45\textwidth}
  \begin{flushleft} %\large
    Vº. Bº. del Tutor:\\[2cm]
    D. José Manuel Aroca Fernández
  \end{flushleft}
\end{minipage}
\hfill
\begin{minipage}{0.45\textwidth}
  \begin{flushleft} %\large
    Vº. Bº. del co-tutor:\\[2cm]
    D. Rodrigo Pascual García
  \end{flushleft}
\end{minipage}
\hfill

\vfill

% para casos con solo un tutor comentar lo anterior
% y descomentar lo siguiente
%Vº. Bº. del Tutor:\\[2cm]
%D. nombre tutor


\newpage\null\thispagestyle{empty}\newpage




\frontmatter

% Abstract en castellano
\renewcommand*\abstractname{Resumen}
\begin{abstract}
La instalación inicial de un equipo obliga habitualmente a localizar y descargar cada aplicación desde una página distinta. Este Trabajo de Fin de Grado presenta \emph{Batch Downloader}, una plataforma web que centraliza un catálogo de software y permite seleccionar varias aplicaciones para obtener sus instaladores oficiales agrupados en un único archivo ZIP.

La plataforma no mantiene ejecutables de forma permanente. Cada solicitud crea un trabajo asíncrono que resuelve y valida fuentes externas, descarga temporalmente los archivos y prepara un ZIP que Core entrega mediante streaming y limpia según el estado de guardado y la actividad. El scraper y la búsqueda semántica constituyen dos partes centrales del proyecto. El sistema admite descargas anónimas, bundles oficiales y personales, acceso de usuarios mediante enlaces de un solo uso por correo y una interfaz en español e inglés.

La búsqueda ofrece modos literal y semántico. El servicio semántico usa un modelo local aprovisionado en la imagen o en el volumen persistente, sin descargar pesos dinámicamente; PostgreSQL con pgvector mantiene una proyección reconstruible, mientras MySQL conserva la autoridad sobre filtros, estados, facetas y paginación. Ante un índice incompleto, una consulta no elegible o un fallo, la petición completa se degrada a modo literal. Esta ampliación no introduce recomendaciones generativas ni creación automática de bundles.

La arquitectura utiliza Vite y TypeScript en el cliente, Java 26 y Spring Boot en el servidor, MySQL para los datos transaccionales, PostgreSQL con pgvector para la proyección semántica, RabbitMQ para el procesamiento asíncrono, Playwright para la resolución dinámica de páginas y contenedores Docker. La evaluación relaciona requisitos, riesgos y pruebas en dos cortes identificados de Git; los resultados no ejecutados se distinguen de los verificados.
\end{abstract}

\renewcommand*\abstractname{Descriptores}
\begin{abstract}
descarga de software, bundles, Spring Boot, Playwright, procesamiento asíncrono, búsqueda semántica, embeddings locales, pgvector, seguridad web
\end{abstract}

\clearpage

% Abstract en inglés
\renewcommand*\abstractname{Abstract}
\begin{abstract}
Setting up a computer usually requires locating and downloading every application from a different website. This Bachelor's Degree Final Project presents \emph{Batch Downloader}, a web platform that provides a centralized software catalogue and allows users to select several applications and obtain their official installers grouped into a single ZIP archive.

The platform does not store executable files permanently. Each request creates an asynchronous job that resolves and validates external sources, temporarily downloads the files, and prepares a ZIP archive. Core streams the archive to the user and cleans up storage according to save and activity state. Source scraping and semantic search are both central to the project. The system supports anonymous downloads, official and user-owned bundles, single-use email links for user sign-in, and a Spanish and English interface.

Search supports lexical and semantic modes. A locally provisioned embedding model runs without dynamic weight downloads; PostgreSQL with pgvector stores a rebuildable projection, while MySQL remains authoritative for filters, status, facets, and pagination. An incomplete or unavailable semantic index, an ineligible query, or a service failure causes the entire request to fall back to lexical search. This extension does not add generative recommendations or automatic bundle creation.

The architecture uses Vite and TypeScript for the client, Java 26 and Spring Boot for the server, MySQL for transactional data, PostgreSQL with pgvector for the semantic projection, RabbitMQ for asynchronous processing, Playwright for dynamic page resolution, and Docker containers. The evaluation links requirements, risks, and tests across two identified Git revisions and distinguishes executed results from unverified claims.
\end{abstract}

\renewcommand*\abstractname{Keywords}
\begin{abstract}
software download, bundles, Spring Boot, Playwright, asynchronous processing, semantic search, local embeddings, pgvector, web security
\end{abstract}

\clearpage

% Indices
\tableofcontents

\clearpage

\listoffigures

\clearpage

\listoftables
\clearpage

\chapter*{Acrónimos}
\addcontentsline{toc}{chapter}{Acrónimos}
\markboth{Acrónimos}{Acrónimos}

Esta lista reúne las abreviaturas utilizadas en la memoria y en la documentación
técnica, junto con su significado.

\begingroup
\small
\setlength{\tabcolsep}{4pt}
\begin{longtable}{p{0.18\textwidth}p{0.71\textwidth}}
\toprule
\textbf{Sigla} & \textbf{Significado}\\
\midrule
\endfirsthead

\multicolumn{2}{l}{\small Acrónimos (continuación)}\\
\toprule
\textbf{Sigla} & \textbf{Significado}\\
\midrule
\endhead

\midrule
\multicolumn{2}{r}{\small Continúa en la página siguiente}\\
\endfoot

\bottomrule
\endlastfoot

AMQP &
\emph{Advanced Message Queuing Protocol} (protocolo avanzado de colas de mensajes)\\

API &
\emph{Application Programming Interface} (interfaz de programación de aplicaciones)\\

CAPTCHA &
\emph{Completely Automated Public Turing test to tell Computers and Humans Apart}\\

CD &
\emph{Continuous Delivery/Deployment} (entrega o despliegue continuos)\\

CDN &
\emph{Content Delivery Network} (red de distribución de contenidos)\\

CI &
\emph{Continuous Integration} (integración continua)\\

CP &
Caso de prueba\\

CRUE &
CRUE Universidades Españolas; originalmente, Conferencia de Rectores de las Universidades Españolas\\

CSRF &
\emph{Cross-Site Request Forgery} (falsificación de petición entre sitios)\\

CSS &
\emph{Cascading Style Sheets} (hojas de estilo en cascada)\\

CU &
Caso de uso\\

DLQ &
\emph{Dead Letter Queue} (cola de mensajes no procesables)\\

DNS &
\emph{Domain Name System} (sistema de nombres de dominio)\\

DTO &
\emph{Data Transfer Object} (objeto de transferencia de datos)\\

E2E &
\emph{End to End} (extremo a extremo)\\

ES &
\emph{ECMAScript}\\

ETag &
Etiqueta HTTP de versión para validación de caché\\

GHCR &
\emph{GitHub Container Registry} (registro de imágenes de contenedor de GitHub)\\

GIL &
\emph{Global Interpreter Lock} (bloqueo global del intérprete de Python)\\

H2 &
Base de datos relacional embebida utilizada por los \emph{inbox} locales\\

HMAC &
\emph{Hash-based Message Authentication Code} (código de autenticación basado en hash)\\

HTML &
\emph{HyperText Markup Language} (lenguaje de marcado de hipertexto)\\

HTTPS &
\emph{HyperText Transfer Protocol Secure}\\

IA &
\emph{Inteligencia artificial}\\

ID &
\emph{Identifier} (identificador)\\

IP &
\emph{Internet Protocol} (protocolo de Internet)\\

ISO &
\emph{International Organization for Standardization} (Organización Internacional de Normalización)\\

JDK &
\emph{Java Development Kit} (kit de desarrollo de Java)\\

JIT &
\emph{Just in Time} (revalidación realizada justo antes de la descarga)\\

JPA &
\emph{Jakarta Persistence API}\\

LLM &
\emph{Large Language Model} (modelo de lenguaje de gran tamaño)\\

LOPDGDD &
Ley Orgánica de Protección de Datos Personales y garantía de los derechos digitales\\

LPI &
Ley de Propiedad Intelectual\\

LSSI &
Ley de servicios de la sociedad de la información y de comercio electrónico\\

LTS &
\emph{Long-Term Support} (soporte a largo plazo)\\

MIME &
\emph{Multipurpose Internet Mail Extensions} (tipo de contenido de Internet)\\

ODS &
\emph{Objetivos de Desarrollo Sostenible}\\

ORM &
\emph{Object-Relational Mapping} (mapeo objeto-relacional)\\

REST &
\emph{Representational State Transfer} (transferencia de estado representacional)\\

RF &
Requisito funcional\\

RGPD &
Reglamento General de Protección de Datos\\

RNF &
Requisito no funcional\\

SMTP &
\emph{Simple Mail Transfer Protocol} (protocolo simple de transferencia de correo)\\

{\footnotesize SOS1--SOS4} &
Competencias de sostenibilidad 1 a 4\\

SSRF &
\emph{Server-Side Request Forgery} (falsificación de peticiones del lado del servidor)\\

SSE &
\emph{Server-Sent Events} (eventos enviados por el servidor)\\

URL &
\emph{Uniform Resource Locator} (localizador uniforme de recursos)\\

UUID &
\emph{Universally Unique Identifier} (identificador universal único)\\

VPS &
\emph{Virtual Private Server} (servidor privado virtual)\\

WCAG &
\emph{Web Content Accessibility Guidelines} (pautas de accesibilidad para el contenido web)\\

ZIP &
Formato de archivo comprimido usado para entregar el resultado de un trabajo\\

\end{longtable}
\endgroup

\clearpage
\chapter*{Glosario}
\addcontentsline{toc}{chapter}{Glosario}
\markboth{Glosario}{Glosario}

Los términos siguientes se utilizan con este significado en la memoria y en los anexos.

\begin{longtable}{p{0.24\textwidth}p{0.65\textwidth}}
\toprule
\textbf{Término} & \textbf{Significado en el proyecto}\\
\midrule
\endfirsthead
\toprule
\textbf{Término} & \textbf{Significado en el proyecto}\\
\midrule
\endhead
\bottomrule
\endlastfoot
Bundle & Colección persistida de aplicaciones, oficial o propiedad de una cuenta; se distingue de la selección temporal de una descarga.\\
Candidato & Enlace o alternativa de instalación detectada para una aplicación antes de su validación final.\\
Enlace mágico & Enlace de acceso enviado por correo, de un solo uso y vigencia limitada; su consumo crea o autentica una cuenta de usuario.\\
Fallback literal & Repetición completa de una búsqueda en modo literal cuando la ruta semántica no puede ofrecer un resultado coherente.\\
Fuente & Página o servicio del editor y configuración asociada desde los que se intenta obtener un instalador.\\
Índice semántico & Proyección reconstruible de documentos y vectores que permite proponer candidatos a una búsqueda; MySQL conserva los filtros y estados públicos.\\
Riesgo residual & Riesgo estimado que permanece después de aplicar una medida; no equivale por sí solo a una prueba superada.\\
Trabajo & Solicitud asíncrona de descarga con estados, elementos y un resultado temporal propio.\\
Trazabilidad & Relación verificable entre requisitos, casos de uso, riesgos, pruebas, resultados y evidencias.\\
\end{longtable}

\clearpage

\mainmatter
\capitulo{1}{Introducción}

\section{Contexto y problema}

La preparación de un ordenador nuevo, la reinstalación del sistema operativo o la configuración de un puesto de trabajo suelen exigir la descarga de numerosas aplicaciones. Aunque cada fabricante publica sus instaladores en páginas oficiales, el usuario debe localizar dichas páginas, identificar la versión adecuada para su sistema y repetir el proceso para cada programa. Esta tarea es mecánica, consume tiempo y aumenta la posibilidad de utilizar enlaces obsoletos o sitios no oficiales.

\emph{Batch Downloader} es una plataforma web que centraliza el descubrimiento y la descarga de software sin convertirse en un repositorio permanente de ejecutables. El catálogo mantiene información descriptiva, metadatos y referencias cifradas a instaladores; los binarios se vuelven a validar y se obtienen cuando el usuario solicita una descarga. El worker reúne temporalmente los archivos aceptados, construye un ZIP con su manifiesto y elimina los artefactos al expirar el resultado.

El proyecto comparte con servicios como Ninite la idea de reducir las acciones necesarias para preparar un equipo. Ninite permite seleccionar varias aplicaciones y generar un instalador personalizado, descarga cada programa desde el sitio de su editor y verifica firmas o resúmenes antes de ejecutarlo~\cite{ninite}. Batch Downloader se diferencia por ofrecer un catálogo administrable y observable, filtros por plataforma, bundles mantenidos por la administración, generación de ZIP en servidor y un pipeline desacoplado que conserva el historial de resolución.

\section{Motivación}

La motivación principal es simplificar un proceso frecuente sin renunciar a la trazabilidad. El usuario debe conocer qué aplicaciones ha elegido, desde qué dominio se obtiene cada instalador y por qué un archivo no ha podido incluirse. Esta transparencia es especialmente importante porque la plataforma procesará ejecutables de terceros y dependerá de páginas externas que pueden cambiar sin previo aviso.

El sistema también aborda un problema de organización. Una selección habitual, como un entorno de desarrollo, un conjunto de utilidades multimedia o las aplicaciones básicas para un equipo nuevo, puede expresarse como un \emph{bundle}. La implementación permite consultar y descargar estas colecciones, mientras que su alta y mantenimiento pertenecen al panel administrativo. Las funciones comunitarias de autoría, clonación, personalización y valoración no forman parte de la entrega.

El buscador ofrece un modo literal y otro híbrido. El primero combina coincidencias normalizadas de nombre, editor y descripción; el segundo fusiona ese orden con la similitud de embeddings locales mediante \emph{Reciprocal Rank Fusion}. En ambos casos MySQL conserva la autoridad sobre estados, filtros, facetas, totales y paginación. PostgreSQL con pgvector mantiene una proyección reconstruible y el sistema degrada la petición completa al modo literal cuando el índice semántico no está disponible o no puede enumerarse de forma exacta.

Los pesos de los modelos se obtienen de Hugging Face, pero las descripciones, los metadatos y las consultas no se envían a ese servicio. Groq o DeepSeek continúan limitados al enriquecimiento opcional de descripciones largas; no participan en el ranking ni crean bundles.

\section{Alcance de esta memoria}

Esta memoria documenta una iteración implementada de Batch Downloader. El repositorio contiene el catálogo web, los servicios Java, el scraper Python, los contratos de mensajería y los flujos de descarga asíncrona. La redacción diferencia los comportamientos verificados mediante pruebas automatizadas de los trabajos que todavía requieren validación de rendimiento o de despliegue público.

El alcance funcional incluye:

\begin{itemize}
	\item consulta y filtrado de un catálogo de aplicaciones.
	\item selección individual o múltiple y generación de un ZIP en servidor.
	\item sesiones para usuarios anónimos, registrados y administradores.
	\item consulta y descarga de bundles, con mantenimiento administrativo.
	\item solicitud de nuevas aplicaciones y envío opcional de avisos por correo.
	\item resolución periódica de fuentes, revalidación inmediata y revisión administrativa.
	\item clasificación persistente en disponible, revisión y sin instalador, con contadores incrementales.
	\item búsqueda literal e híbrida con selector persistente, fallback íntegro y versiones explícitas de modelo e índice.
	\item observabilidad de trabajos, colas, reintentos y cambios del catálogo.
\end{itemize}


La entrega no declara una medición de capacidad en producción. Sí incorpora manuales de programación y de usuario contrastados con el repositorio, los contratos OpenAPI y AsyncAPI, pruebas automatizadas para los cuatro servicios Java, el scraper y el frontend, y una verificación visual del flujo de catálogo y descarga.

\section{Estructura del documento}

\begin{itemize}
	\item El capítulo 2 define los objetivos.
	\item El capítulo 3 presenta los conceptos necesarios para comprender la solución.
	\item El capítulo 4 justifica las tecnologías y técnicas seleccionadas.
	\item El capítulo 5 describe la arquitectura y los aspectos más relevantes del diseño.
	\item El capítulo 6 compara la propuesta con Ninite, Home Updater y UBULabs.
	\item El capítulo 7 resume las conclusiones alcanzadas durante el análisis y plantea líneas de trabajo futuras.
	\item La documentación técnica se entrega por separado. Sus anexos contienen el plan de proyecto, la especificación de requisitos, el diseño de datos y procedimientos, los manuales de programación y de usuario, y el análisis de sostenibilidad.
\end{itemize}

\capitulo{2}{Objetivos del proyecto}

\section{Objetivo general}

El objetivo general es diseñar e implementar una plataforma web que permita localizar aplicaciones, seleccionar uno o varios instaladores y descargarlos agrupados desde sus fuentes oficiales, evitando el almacenamiento permanente de ejecutables y ofreciendo información clara sobre el origen y el estado de cada archivo.

\section{Objetivos funcionales}

\begin{enumerate}
	\item Construir un catálogo administrable de software con nombre, descripción, etiquetas, página oficial, sistemas operativos, arquitecturas y estado de disponibilidad.
	\item Permitir búsqueda por nombre, editor o descripción y combinarla con filtros de estado, sistema operativo, etiquetas y editor, manteniendo la paginación en el servidor.
	\item Permitir la selección de una o varias aplicaciones. Una selección múltiple se tratará como un bundle temporal para generar un ZIP.
	\item Consultar y descargar bundles oficiales y propios mediante enlaces estables; cada cuenta puede crear, editar y eliminar sus bundles públicos o privados desde el panel personal.
	\item Proporcionar panel personal, historial de descargas, verificación de correo y recuperación de cuenta.
	\item Proporcionar al administrador funciones para mantener aplicaciones, etiquetas, reglas de resolución, evidencias de fuente y estados del catálogo.
	\item Clasificar cada aplicación activa de forma exhaustiva como disponible, en revisión o sin instalador, manteniendo contadores incrementales coherentes con el catálogo.
	\item Informar del progreso de cada descarga, de los archivos incluidos y de los errores que impidan completar una parte de la solicitud.
\end{enumerate}

\section{Objetivos de descarga y procesamiento}

\begin{enumerate}
	\item Resolver dinámicamente el enlace del instalador desde una fuente oficial mediante una estrategia específica o, como respaldo, mediante navegación automatizada con Playwright.
	\item Validar protocolo, dominio final, redirecciones, tipo de contenido, extensión y tamaño antes de aceptar un archivo.
	\item Descargar los instaladores en el servidor, generar un ZIP con un manifiesto, incluir accesos \ruta{.url} cuando solo exista una página oficial y eliminar todos los archivos temporales al terminar, cancelar o expirar el trabajo.
	\item Procesar las solicitudes de manera asíncrona con RabbitMQ para separar la interacción web de las operaciones de resolución, descarga y compresión.
	\item Aplicar reintentos controlados y permitir resultados \emph{PARTIAL} o \emph{MANUAL\_ONLY} cuando una aplicación falle o solo admita instalación manual, manteniendo una explicación individual para cada incidencia.
	\item Calcular, cuando las fuentes lo permitan, un tamaño estimado antes de iniciar el trabajo y aplicar límites por archivo, ZIP, usuario, dirección IP y número de operaciones concurrentes.
\end{enumerate}

\section{Estado de consecución}

La iteración documentada materializa el flujo esencial de catálogo y descarga. El catálogo devuelve los sistemas operativos con un candidato estructuralmente válido y el cliente permite combinarlos con semántica OR mediante parámetros repetidos en la URL. Una comprobación antigua no retira por sí sola una aplicación: la fuente se revalida justo antes de la descarga y solo un resultado terminal explícito modifica su clasificación. MySQL materializa el estado público y sus cuatro contadores mediante columnas generadas y triggers transaccionales.

La creación de un trabajo es asíncrona: el Core persiste el trabajo, sus elementos y el evento de salida en una única transacción; RabbitMQ entrega solo identificadores al worker; y el resultado se publica como progreso, estado terminal y URL temporal de MinIO.

La descarga funciona tanto para selecciones como para bundles. El servidor resuelve un bundle completo, no la previsualización mostrada por la interfaz, impone el máximo de cien aplicaciones y comunica las aplicaciones solicitadas, aceptadas y omitidas. Los trabajos anónimos se asocian a una cookie aleatoria cuyo HMAC se almacena en la base de datos; no se incluyen tokens de propiedad en mensajes, URL ni registros. Los trabajos autenticados y anónimos comparten los estados \emph{QUEUED}, \emph{RESOLVING}, \emph{DOWNLOADING}, \emph{PACKAGING}, \emph{READY}, \emph{PARTIAL}, \emph{MANUAL\_ONLY}, \emph{FAILED}, \emph{CANCELLED} y \emph{EXPIRED}. Para Linux se seleccionan perfiles APT, DNF, Pacman o Zypper y se conservan las comprobaciones SHA-256/GPG y la posibilidad de rollback.

\section{Objetivos de búsqueda semántica}

\begin{enumerate}
	\item Construir un documento canónico por aplicación con nombre, package ID, editor, etiquetas, descripciones, plataformas, arquitecturas, versión y dominio registrado, acompañado de un hash determinista.
	\item Representar esos documentos mediante un modelo de embeddings local aprovisionado y mantenerlos en PostgreSQL con pgvector como una proyección reconstruible del catálogo MySQL.
	\item Interpretar consultas en lenguaje natural, enumerar candidatos mediante similitud coseno y aplicar después los filtros ordinarios en MySQL.
	\item Validar la revisión, dimensión, manifiesto y cobertura del modelo local sin descargar pesos dinámicamente.
	\item Degradar de forma íntegra al modo literal ante timeout, error, falta de cobertura o consulta no elegible; el índice vectorial nunca decide estados, totales o permisos.
\end{enumerate}

\section{Objetivos técnicos y de calidad}

\begin{itemize}
	\item Desarrollar un cliente con React, React Router, Vite, HTML, CSS y TypeScript, con las versiones de Vite/Vitest fijadas por el lockfile.
	\item Ejecutar los servicios Java con JDK 26 y generar bytecode Java 25, manteniendo una separación explícita entre Core API, worker de descarga, notificaciones y traducciones.
	\item Documentar la API HTTP mediante OpenAPI, cuyo propósito es describir de forma independiente del lenguaje las capacidades de una API~\cite{openapi}.
	\item Utilizar MySQL para la información transaccional y su proyección incremental del catálogo; mantener PostgreSQL con pgvector como proyección semántica reconstruible.
	\item Ejecutar el scraper sobre CPython 3.14.6 \emph{free-threaded}, comprobar que el GIL continúa deshabilitado tras importar extensiones y reservar los hilos para funciones puras de CPU.
	\item Contenerizar los componentes con Docker y desplegarlos en Compose local y en el VPS de Oracle mediante Coolify, usando imágenes multi-arquitectura de GHCR y tags inmutables.
	\item Configurar Resend/SMTP, healthchecks, métricas expuestas, copias OCI y monitorización diaria sin afirmar un servidor Prometheus central.
	\item Definir pruebas unitarias, de integración, de seguridad, de resolución de enlaces, de generación de ZIP y de extremo a extremo.
	\item Mantener trazabilidad mediante registros de administración, resolución, descargas y uso del asistente.
\end{itemize}

\section{Objetivos de seguridad y sostenibilidad}

El sistema reduce los riesgos derivados de procesar URL externas y ejecutables validando HTTPS, host, puerto, todas las direcciones DNS y la IP efectiva tras cada redirección. Se rechazan redes privadas, reservadas, de enlace local y puntos de metadatos; no se degrada a HTTP. A ello se añaden sanitización de nombres, validación de MIME y firma, control de acceso, CSRF y cuotas de uso. La plataforma no ejecuta los instaladores descargados.

Desde el punto de vista de sostenibilidad, se buscará evitar copias permanentes, eliminar temporales, limitar trabajos innecesarios, reutilizar metadatos válidos y escalar los servicios según su carga. También se atenderá a la accesibilidad, la claridad de la información, la protección de datos y el acceso equitativo a la funcionalidad básica.

\capitulo{3}{Conceptos teóricos}

\section{Catálogo centralizado y fuentes oficiales}

Una plataforma de descarga puede almacenar los binarios o limitarse a mantener información para encontrarlos. Batch Downloader adopta el segundo enfoque: conserva metadatos y reglas de resolución, mientras que los instaladores se obtienen bajo demanda. Esta separación reduce la duplicación de ejecutables y evita que el catálogo se convierta en una fuente desactualizada respecto al editor original.

El concepto de \emph{fuente oficial} no se limita a la URL visible de una página. Una descarga puede atravesar redirecciones o terminar en una red de distribución de contenidos. La URL final se comprueba antes de aceptar el archivo: se exige HTTPS, se inspeccionan host, puerto, respuestas DNS e IP efectiva y se rechazan destinos internos o reservados. La interfaz muestra la página oficial, el estado y la fecha de la última comprobación disponible.

\section{Bundles administrados}

Un bundle es una colección identificable de aplicaciones. Puede representar una necesidad concreta, por ejemplo ``herramientas para desarrollo web'', y proporciona un enlace estable para consultar o descargar la selección completa. En la entrega se distinguen:

\begin{description}
	\item[Bundle oficial:] colección creada y mantenida desde la administración, consultable mediante la API pública cuando su estado lo permite.
	\item[Bundle propio:] colección creada por una cuenta autenticada. Puede ser privada, accesible solo desde el panel personal, o pública, en cuyo caso aparece en el catálogo de bundles.
	\item[Selección temporal:] conjunto de aplicaciones elegido directamente en el catálogo para una descarga concreta.
\end{description}

La API pública permite listar, consultar y descargar bundles oficiales y propios públicos. Las rutas bajo \ruta{/api/v1/users/me/bundles} permiten crear, editar y eliminar los bundles de la cuenta autenticada; el control de propiedad usa el UUID de la cuenta y una versión esperada para detectar ediciones concurrentes. El panel personal también reúne el historial de descargas. No se presenta como implementada la valoración comunitaria ni la clonación.

\section{Resolución de enlaces y scraping}

Las URL de descarga cambian por versiones, regiones, sistemas operativos o decisiones del proveedor. Un \emph{resolver} es el componente que transforma la información de una aplicación y el contexto solicitado en una URL descargable validada.

La estrategia implementada es progresiva y combina información estructurada con navegación controlada:

\begin{enumerate}
	\item consultar las API y datos JSON de Winstall, incluido \ruta{\_\_NEXT\_DATA\_\_};
	\item priorizar la web oficial y resolvers específicos, como GitHub Releases;
	\item analizar HTML y candidatos HTTP con reglas acotadas;
	\item navegar por la página oficial mediante Playwright como último respaldo, inspeccionando enlaces, botones, redirecciones y tráfico de red;
	\item marcar la fuente para revisión manual si no se obtiene un resultado seguro.
\end{enumerate}

Playwright permite controlar navegadores modernos en modo sin interfaz~\cite{playwright}. En el proyecto se encapsula como respaldo de los resolvers de Python; no sustituye a las comprobaciones posteriores. El resultado debe mantener HTTPS, resolver a una red pública, cumplir los límites establecidos y superar la validación del binario.

\section{Procesamiento asíncrono}

Resolver varias páginas, descargar archivos grandes y comprimirlos puede superar el tiempo razonable de una petición HTTP. El procesamiento asíncrono separa la creación de la solicitud de su ejecución.

El Core valida la selección y publica un mensaje mediante un patrón outbox en RabbitMQ. El mensaje transporta exclusivamente identificadores. El worker consume el comando de forma idempotente, solicita al scraper la resolución interna autenticada y actualiza el estado a través de eventos. Spring AMQP proporciona las abstracciones de envío y consumo~\cite{springamqp}. El usuario recibe progreso mediante SSE, dispone de sondeo como respaldo y descarga el ZIP cuando el estado es \emph{READY} o \emph{PARTIAL}.

Los estados previstos para un trabajo son:

\begin{center}
	\begin{tabularx}{0.92\textwidth}{p{0.25\textwidth}X}
		\toprule
		\textbf{Estado}    & \textbf{Significado}                                 \\
		\midrule
		\emph{QUEUED}      & Espera de un worker disponible.                      \\
		\emph{RESOLVING}   & Resolución interna de fuentes.                       \\
		\emph{DOWNLOADING} & Descarga temporal y validación de instaladores.      \\
		\emph{PACKAGING}   & Construcción del ZIP y del manifiesto.               \\
		\emph{READY}       & Archivo completo preparado para su entrega.          \\
		\emph{PARTIAL}     & Archivo preparado con al menos un instalador válido. \\
		\emph{FAILED}      & No se ha podido producir un resultado válido.        \\
		\emph{CANCELLED}   & Trabajo cancelado de forma idempotente.              \\
		\emph{EXPIRED}     & Resultado y objetos temporales eliminados.           \\
		\bottomrule
	\end{tabularx}
\end{center}

\section{Generación temporal del ZIP}

Cada trabajo utiliza un espacio aislado identificado por su código. Los instaladores no se incorporan a un almacenamiento permanente del catálogo. MinIO almacena temporalmente staging, el manifiesto y el ZIP para compartir artefactos entre workers, aplicar tiempos de expiración y evitar que el servidor web mantenga archivos grandes en su sistema local.

El ZIP incluye los instaladores disponibles en su raíz y un fichero \ruta{manifest.json}. El manifiesto registra la aplicación, el sistema operativo, la arquitectura, el nombre del archivo y el resultado del elemento. Si alguna descarga falla, el trabajo puede terminar como \emph{PARTIAL} y conservar el ZIP con los archivos restantes; si no se obtiene ninguno, finaliza como \emph{FAILED}.

\section{Embeddings y búsqueda semántica}

Un embedding es una representación numérica de un elemento en un espacio vectorial. Textos con significado relacionado tienden a ocupar posiciones próximas según una medida de distancia.

PostgreSQL con pgvector permite almacenar vectores y realizar búsquedas de vecinos cercanos mediante distancia euclídea, producto interno o similitud coseno~\cite{pgvector}.

En Batch Downloader el vector no sustituye al registro transaccional. El scraper produce un documento canónico limitado a metadatos públicos y calcula un \ruta{contentHash}; el indexador vuelve a generar el embedding cuando cambia ese hash y marca como inactivos los documentos que desaparecen del barrido. El modelo activo se aprovisiona en la imagen o en un volumen persistente con un manifiesto de revisión; el servicio funciona con \ruta{local\_files\_only} y no descarga pesos durante el arranque.

La búsqueda semántica enumera candidatos mediante similitud coseno y entrega a Core una colección acotada. Si falta el modelo, el índice está incompleto, la consulta no cumple el mínimo o la respuesta no puede enumerarse de forma íntegra, Core descarta la colección y repite la petición completa en modo literal. MySQL aplica después los filtros estructurales, los estados, las facetas, los desempates y la paginación.

No se afirma una fusión RRF ni una familia de modelos entrenables: la degradación segura a búsqueda literal es la garantía funcional verificable de esta versión.

\section{Modelo de lenguaje para enriquecimiento de descripciones}

El modelo de lenguaje recibe metadatos acotados de una aplicación para producir una descripción larga. Esta tarea se ejecuta en una cola persistente, fuera del descubrimiento y de la descarga, y su resultado nunca decide si una fuente es segura.

El enriquecimiento admite Groq y DeepSeek mediante una interfaz desacoplada~\cite{groq}. La configuración selecciona el proveedor y los modelos sin modificar la cola ni el catálogo; los límites de concurrencia y cuota impiden que este proceso bloquee el pipeline principal.

Las respuestas se validan antes de persistirse. Las responsabilidades implementadas del modelo quedan limitadas a:

\begin{itemize}
	\item redactar una descripción larga a partir de metadatos existentes;
	\item registrar proveedor, modelo, fecha y error de generación;
	\item respetar una descripción manual y permitir reintentos controlados.
\end{itemize}

El modelo de lenguaje no puede aprobar dominios, resolver una URL, ejecutar descargas, cambiar reglas de scraping ni generar un ZIP. La búsqueda semántica usa modelos de embeddings locales y no incorpora recomendaciones generativas, explicaciones automáticas ni creación asistida de bundles.

\section{Seguridad en la descarga de recursos externos}

El backend recibirá y seguirá URL externas, lo que introduce riesgo de falsificación de peticiones del lado del servidor o SSRF. Un atacante podría intentar acceder a servicios internos, metadatos de infraestructura o redes privadas. La mitigación prevista incluye:

\begin{itemize}
	\item admitir exclusivamente HTTPS;
	\item resolver DNS y rechazar direcciones locales, privadas, de enlace local o reservadas;
	\item repetir la validación después de cada redirección;
	\item no degradar a HTTP y comprobar la red pública después de cada redirección;
	\item limitar redirecciones, tiempo, tamaño y concurrencia;
	\item no ejecutar ningún archivo y sanitizar su nombre antes de incorporarlo al ZIP.
\end{itemize}

También se aplicarán controles de autenticación, autorización, limitación de frecuencia, cuotas y auditoría. Estas medidas evitan que la plataforma se utilice como proxy de descargas masivas y proporcionan evidencias para investigar incidencias.

\capitulo{4}{Técnicas y herramientas}

\section{Criterios de selección}

Las tecnologías se han seleccionado por su tipado, separación de responsabilidades, facilidad de despliegue, observabilidad y posibilidad de escalar los procesos costosos de forma independiente. La selección se ha contrastado mediante la compilación de los servicios Java, las pruebas del scraper y del frontend, la validación de los contratos públicos y los smoke tests del despliegue en Oracle/Coolify.

\section{Mapa tecnológico}

\begin{landscape}
\begin{figure}[p]
\centering
\includegraphics[width=0.96\linewidth,height=0.90\textheight,keepaspectratio]{system_architecture_layers.pdf}
\caption{Arquitectura de sistema de Batch Downloader.}
\label{fig:arquitectura-capas}
\end{figure}
\end{landscape}
\FloatBarrier

\begin{figure}[p]
\centering
\includegraphics[width=\linewidth,height=0.90\textheight,keepaspectratio]{technology_stack.pdf}
\caption{Tecnologías del proyecto.}
\label{fig:mapa-tecnologico}
\end{figure}
\FloatBarrier

\section{Frontend}

\subsection{React, React Router, TypeScript y Vite}

El cliente utiliza React~18.3 y React DOM para componer la interfaz, React Router~7.18 para representar navegación y filtros en la URL, y TypeScript~5.9 para tipar contratos, estados y formularios. Vite~8.2 proporciona el servidor de desarrollo y la construcción optimizada~\cite{vite}; el plugin oficial de React integra la transformación JSX. Lucide React aporta iconos de interfaz sin introducir una biblioteca de marcas.

HTML y CSS mantienen una tabla compacta y adaptable. Los eventos SSE notifican el progreso de trabajos y WebSocket anuncia cambios del catálogo y de la administración. \emph{translation-service} sirve catálogos JSON en español e inglés; el cliente detecta el idioma del navegador y permite cambiarlo en cualquier momento.

\subsection{Entrega y pruebas del cliente}

La imagen de producción se construye con Node.js~24 y se entrega mediante Nginx~1.27, que sirve los recursos estáticos y enruta \ruta{/api} al Core. Vitest~4.1, Testing Library, jest-dom y jsdom verifican componentes y flujos sin navegador completo; el proveedor V8 produce cobertura. Todas estas pruebas y código no se incluyen en la imagen de producción.

\section{Backend}

\subsection{JDK 26 y bytecode Java 25}

Java es el lenguaje de los servicios del servidor por su sistema de tipos, madurez y bibliotecas. Los contenedores se ejecutan con JDK 26~\cite{java26}, mientras que Maven genera bytecode Java 25. Con Java 26 me permito tener la ventaja de correciones, warnings y rendimiento que me ofrece el lenguaje dentro del desarrollo, incluyendo las respectivas extensiones de VS Code, mientras que con Java 25 me aseguro de que el bytecode generado sea compatible con todas las librerías y dependencias como, por ejemplo, SonarQube.
\subsection{Spring Boot}

Los cuatro servicios Java usan Spring Boot~3.5.16~\cite{springboot}. Sus dependencias directas se distribuyen por responsabilidad:

\begin{description}
	\item[Spring Web MVC y WebSocket:] API REST, SSE, comprobaciones de salud y notificación de cambios del catálogo.
	\item[Spring Security, BCrypt y Spring Session JDBC:] acceso de usuarios mediante enlaces de un solo uso, credencial administrativa con hash, sesiones persistentes, CSRF, cookies seguras, roles y protección de rutas~\cite{springsession}.
	\item[Spring Data JPA y Spring JDBC:] repositorios transaccionales para identidad, outbox y trabajos, junto con consultas explícitas para las proyecciones de catálogo~\cite{springdatajpa}.
	\item[Spring Validation:] validación declarativa de entradas y configuraciones.
	\item[Spring Actuator:] salud, métricas y observabilidad de servicios.
	\item[Spring AMQP:] publicación y consumo de comandos y eventos mediante RabbitMQ.
	\item[Spring Mail:] envío SMTP desde el servicio de notificaciones.
\end{description}

Flyway administra las tablas propiedad del Core; el scraper mantiene las suyas con Alembic. El driver MySQL conecta Core con la base compartida, el SDK de MinIO gestiona objetos y H2 conserva los inbox locales del worker y de notificaciones. \emph{springdoc-openapi}~2.6 expone documentación navegable del contrato.

Las pruebas Java usan JUnit~5, Spring Test, Mockito, AssertJ y Spring Rabbit Test. Testcontainers ejecuta los casos que requieren semántica real de MySQL~8.4~\cite{testcontainers}; H2 se limita a pruebas y a los inbox locales que no dependen de funciones MySQL. JaCoCo~0.8.14 genera cobertura. ArchUnit verifica la arquitectura y las dependencias de paquetes~\cite{archunit}.

\subsection{OpenAPI}

OpenAPI describe el contrato HTTP versionado bajo \ruta{/api/v1}~\cite{openapi}. AsyncAPI documenta los intercambios RabbitMQ y JSON Schema valida cada evento versionado. Estos contratos se comprueban en integración continua y no se utilizan para realizar scraping.

\section{Resolución de fuentes}

\subsection{Playwright como respaldo del scraper Python}

Playwright controla un navegador en modo \emph{headless} para páginas que requieren JavaScript, interacción o inspección de red~\cite{playwright}. En Batch Downloader se ejecuta desde el scraper Python como último recurso del registro de estrategias, después de los resolvers conocidos, GitHub Releases y análisis HTML. Cada proveedor limita la automatización a pasos conocidos y selectores revisables.

La navegación generalista se mantendrá como mecanismo de respaldo. Los proveedores con procesos complejos podrán disponer de resolvers específicos, de forma que los cambios en una página no afecten al resto del catálogo.

\subsection{Scraper y supervisor de enriquecimiento en Python}

La imagen de producción compila CPython~3.14.6 desde la fuente verificada con \ruta{--disable-gil}. Al arrancar comprueba que el GIL permanece deshabilitado después de importar las extensiones nativas; una variable de entorno no puede ocultar una incompatibilidad. FastAPI y Uvicorn proporcionan la API y el servidor ASGI~\cite{fastapi}; Pydantic Settings valida la configuración. SQLAlchemy, Alembic y aiomysql gestionan el catálogo y sus migraciones MySQL~\cite{sqlalchemy,alembic}; aiosqlite se utiliza en pruebas ligeras. La distribución fuente de SQLAlchemy se instala con \ruta{DISABLE\_SQLALCHEMY\_CEXT=1}, tanto en 3.14t como en el control estándar, para usar su implementación Python y evitar que una extensión C reactive el GIL.

HTTPX con HTTP/2 realiza peticiones controladas; selectolax analiza HTML y Playwright~1.61 resuelve páginas que requieren JavaScript. APScheduler inicia ejecuciones periódicas, Tenacity aplica reintentos, tldextract y dnspython apoyan las comprobaciones de dominio y DNS, cryptography cifra las URL, packaging interpreta versiones y structlog produce registros estructurados. Winstall se consulta mediante su API, \ruta{\_\_NEXT\_DATA\_\_} y HTML; GitHub Releases y los recursos de repositorio aportan versiones e iconos cuando existe evidencia suficiente.

Se conserva un único event loop. Un \ruta{ThreadPoolExecutor} acotado recibe exclusivamente parseo HTML, normalización, scoring y deduplicación puros. Las sesiones SQLAlchemy, clientes HTTP, objetos Playwright, eventos, locks y cachés mutables no cruzan hilos. La imagen de producción se valida con CPython 3.14.6t y la configuración de ejecución queda fijada por los manifiestos.

El scheduler adquiere leases y ejecuta \ruta{searcher\_filter}, \ruta{filter\_scraper}, \ruta{scraper\_so\_filter} y \ruta{so\_filter\_descriptor}. El enriquecimiento de descripciones tiene su propia cola persistente; los iconos se resuelven dentro del trabajador de plataforma. pytest, pytest-asyncio y respx cubren el comportamiento asíncrono y HTTP; Ruff y mypy se emplean como comprobaciones estáticas. \ruta{requirements.txt} bloquea las versiones de dependencias de ejecución y prueba para que la imagen y CI sean reproducibles.

\section{Persistencia}

\subsection{MySQL}

MySQL~8.4 almacena los datos transaccionales de identidad, catálogo, fuentes, bundles, trabajos, outbox y pipeline. Core y scraper comparten una instancia, pero cada uno es propietario de grupos de tablas distintos; las solicitudes públicas de aplicaciones ya no forman parte del esquema. La clasificación pública del catálogo se proyecta de forma incremental: los triggers trasladan los cambios de candidatos a fuentes, aplicaciones y una única fila de contadores, de modo que el endpoint de estadísticas no agrega las tablas operativas.

\subsection{H2}

\emph{download-worker} y \emph{notification-service} utilizan bases H2 locales para sus inbox idempotentes. Esta persistencia no sustituye MySQL ni representa una copia del catálogo; evita reprocesar mensajes después de un reinicio. H2 también sirve en pruebas Java que no dependen de funciones específicas de MySQL.

\subsection{PostgreSQL, pgvector y servicio semántico}

La composición incluye PostgreSQL~16 con pgvector y un servicio Python independiente. Sus migraciones crean documentos canónicos, una única revisión local del modelo, embeddings de dimensión validada, trabajos con lease y estado de índice. La comprobación en PostgreSQL rechaza cualquier vector cuya dimensión no coincida con el manifiesto del modelo.

El servicio expone una búsqueda interna autenticada y devuelve candidatos, versión de modelo, versión de índice y una señal de degradación. El indexador reconcilia periódicamente el feed del scraper, actualiza solo hashes modificados y conserva el estado suficiente para reconstruir la proyección. MySQL continúa siendo la fuente de verdad efectiva para los datos públicos.

\subsection{Modelo local y degradación segura}

La imagen del servicio semántico contiene el runtime y recibe el modelo aprobado mediante un volumen persistente identificado por un manifiesto. El arranque valida la revisión, dimensión y hash; no descarga pesos de forma dinámica ni envía textos o consultas a Hugging Face. El indexador solo publica una versión cuando la cobertura del feed es coherente con el \ruta{contentHash} de los documentos activos.

La API conserva dos modos explícitos: \ruta{lexical} y \ruta{semantic}. Si la consulta es demasiado corta, el modelo no está disponible, el índice está incompleto o se produce un error, Core ejecuta de nuevo la petición completa en modo léxico. Así se mantienen deterministas los filtros, las facetas y la paginación sin afirmar benchmarks históricos ni entrenamiento de varios modelos.

El control de coherencia compara la cobertura declarada con el feed vigente y conserva una versión anterior para reconstrucción. No se presenta esta comprobación como un benchmark de modelos ni como entrenamiento de una familia de embeddings.

\section{Mensajería y almacenamiento temporal}

\subsection{RabbitMQ}

RabbitMQ~3.13 desacopla la API de los workers encargados de descargar, comprimir y notificar~\cite{rabbitmq}. Spring AMQP ofrece una abstracción de alto nivel para consumidores asíncronos~\cite{springamqp}. Los exchanges versionados transportan comandos y eventos con \emph{eventId}, tipo, versión de esquema y correlación; la entrega es al menos una vez y los inbox proporcionan idempotencia.

Los mensajes contendrán identificadores, no datos sensibles ni archivos. El estado persistente del trabajo permitirá reintentos controlados y evitará depender exclusivamente de la cola.

\subsection{MinIO}

MinIO almacena temporalmente los ZIP preparados y los objetos intermedios que deben compartirse entre instancias~\cite{minio}. El cliente \ruta{mc} configura el bucket al arrancar. Core conserva el inventario y las reservas de almacenamiento, entrega el ZIP mediante streaming con soporte de rangos y coordina la limpieza al confirmarse el guardado, cancelarse el trabajo o agotarse la actividad. No se utiliza como repositorio permanente de instaladores.

\subsection{SMTP y Mailpit}

El servicio de notificaciones consume eventos, conserva su inbox en H2 y envía correo mediante Spring Mail y SMTP~\cite{smtp}. Mailpit~1.30 captura los mensajes en desarrollo; Resend mediante SMTP entrega en producción el enlace de acceso de un solo uso. Las credenciales administrativas siguen un flujo separado. Las variables de correo se inyectan como secretos y no se escriben en la memoria.

\section{Enriquecimiento mediante LLM}

Groq y DeepSeek se utilizan exclusivamente para generar descripciones largas. El scraper no expone directamente el proveedor al resto de módulos: una interfaz de generación y una estrategia de fallback permiten cambiar modelo, proveedor o límites sin modificar la cola. Se registran el proveedor, modelo, hash de entrada, fecha y error; una descripción manual no se sustituye automáticamente. Los embeddings se calculan localmente en el servicio semántico y no reutilizan estos proveedores.

\section{Contenedores y despliegue}

\subsection{Docker}

Cada componente se empaqueta en una imagen Docker~\cite{docker}. \ruta{docker-compose.yml} construye el árbol local y \ruta{docker-compose.ghcr.yml} consume imágenes publicadas en GHCR. La composición incluye Nginx, Core, scraper API y scheduler, worker, notificaciones, traducciones, servicio e indexador semánticos, MySQL, PostgreSQL/pgvector, RabbitMQ, MinIO/\ruta{mc} y Mailpit. Los healthchecks ordenan migraciones y arranque sin mezclar la propiedad de esquemas ni confundir el puerto publicado de PostgreSQL con su puerto interno 5432.

\section{Herramientas de desarrollo y documentación}

Git mantiene el control de versiones. Maven~3.9.16 compila y verifica el reactor Java; npm usa el lockfile del frontend. GitHub Actions valida contratos, pruebas y construcción, publica imágenes en GHCR y ejecuta el análisis configurado para SonarQube~\cite{sonarqube}.

La memoria se mantiene con \LaTeX{}, pdf\LaTeX{}, latexmk y BibTeX. Mermaid CLI~11.12 es la fuente editable de los diagramas; Puppeteer utiliza Chromium para renderizarlos en SVG y PDF. Los scripts PowerShell reproducen el proceso y Poppler aporta \ruta{pdfinfo}, extracción de texto y conversión a PNG para la inspección visual final.

\section{Resumen de correspondencia}

\begin{table}[H]
	\centering
	\small
	\begin{tabularx}{\textwidth}{p{0.25\textwidth}X}
		\toprule
		\textbf{Necesidad}      & \textbf{Tecnología empleada}                                                 \\
		\midrule
		Cliente web             & React, React Router, TypeScript, Vite, Lucide y Nginx                        \\
		API y lógica de negocio & JDK 26, bytecode Java 25 y Spring Boot                                       \\
		Contrato HTTP y eventos & OpenAPI, AsyncAPI y JSON Schema                                              \\
		Datos transaccionales   & MySQL 8.4, Flyway, Alembic y proyecciones mediante triggers                  \\
		Inbox idempotente       & H2 en worker y notificaciones                                                \\
		Trabajos asíncronos     & RabbitMQ y Spring AMQP                                                       \\
		Scraping dinámico       & CPython 3.14.6 \emph{free-threaded}, FastAPI, HTTPX, selectolax y Playwright \\
		Objetos temporales      & MinIO con expiración                                                         \\
		Descripciones largas    & Groq/DeepSeek mediante una interfaz desacoplada                              \\
		Búsqueda semántica      & PostgreSQL 16, pgvector, runtime de embeddings y modelo local aprovisionado  \\
		Pruebas                 & JUnit, Testcontainers, pytest, Vitest, ArchUnit y Testing Library                      \\
		Portabilidad y entrega  & Docker Compose, GitHub Actions y GHCR                                        \\
		Documentación           & \LaTeX{}, BibTeX, Mermaid, PowerShell y Poppler                              \\
		\bottomrule
	\end{tabularx}
	\caption{Tecnologías empleadas por responsabilidad.}
	\label{tabla:tecnologias-responsabilidad}
\end{table}

\capitulo{5}{Aspectos relevantes del desarrollo del proyecto}

\section{Enfoque de desarrollo}

Las partes más importantes del proyecto se dividen en la creación y correcta funcionalidad del scraper y la IA semántica. Con respecto a la seguridad, el proyecto pide a los usuarios registrarse solo con el correo, de tal forma que no se almacenan contraseñas, no hace falta cambiarla, olvidarse de ella es imposible. Solo necesita el correo. Es el mismo sistema que siguen empresas a día de hoy, como Antrophic (Claude) o Groq INC (Groq).

Las iteraciones materializadas hasta esta entrega son:

\begin{enumerate}
	\item Catálogo, roles, sesión persistente y búsqueda convencional.
	\item Selección de aplicaciones, filtrado por sistema operativo y generación asíncrona de ZIP.
	\item Consulta y descarga de bundles oficiales y públicos o privados, con panel personal e historial de descargas.
	\item Resolvers específicos, colas persistentes, enriquecimiento y observabilidad.
	\item Clasificación materializada, contadores incrementales y revalidación justo antes de descargar.
	\item Búsqueda literal y semántica, proyección pgvector, modelo local aprovisionado y degradación íntegra.
	\item Contratos OpenAPI/AsyncAPI, imágenes de contenedor y validaciones automatizadas.
\end{enumerate}

Esta secuencia separa la generación de descripciones mediante LLM del núcleo de descarga. Se completaron catálogos en español e inglés, se introdujo la entrega observable del ZIP y se amplió la operación del scraper. La evidencia ejecutada y sus límites se registran en el anexo de pruebas. La integración continua configura SonarQube Community Build para el código nuevo~\cite{sonarqube}.

\section{Actores y permisos}

\begin{table}[H]
	\centering
	\small
	\begin{tabularx}{\textwidth}{p{0.19\textwidth}XXX}
		\toprule
		\textbf{Acción}               & \textbf{Anónimo} & \textbf{Registrado}     & \textbf{Administrador} \\
		\midrule
		Consultar catálogo            & Sí               & Sí                      & Sí                     \\
		Generar ZIP                   & Sí, con cuota global   & Sí, con cuotas propias  & Sí                     \\
		Consultar y descargar bundles & Sí               & Sí                      & Sí                     \\
		Mantener bundles oficiales    & No               & No                      & Sí                     \\
		Mantener bundles propios      & No               & Sí 					   & No                     \\
		Mantener catálogo y resolvers & No               & No                      & Sí                     \\
		Revisar registros             & No               & Solo los propios        & Sí                     \\
		\bottomrule
	\end{tabularx}
	\caption{Permisos generales por actor.}
	\label{tabla:permisos-actores}
\end{table}

\section{Arquitectura implementada}

La solución utiliza una arquitectura de servicios con responsabilidades delimitadas:

\begin{description}
	\item[Frontend/Nginx:] presenta catálogo, selección, bundles, panel personal, historial, progreso SSE, administración y traducciones en español e inglés, y enruta la API pública.
	\item[Core API:] valida entradas, aplica permisos, mantiene identidad, bundles y trabajos, y publica eventos mediante outbox.
	\item[Scraper API y scheduler:] gestionan proveedor, fuentes, resoluciones, ejecuciones, estrategias y enriquecimiento.
	\item[Download worker:] consume trabajos idempotentes, descarga y valida archivos, genera ZIP y publica resultados.
	\item[Notification service:] consume eventos de identidad y de trabajos, registra intentos y entrega correo SMTP sin cambiar el resultado del job.
	\item[Translation service:] entrega los catálogos español e inglés versionados con ETag; el frontend conserva una alternativa local para el idioma de referencia.
	\item[Semantic service e indexer:] mantienen la proyección PostgreSQL/pgvector, generan embeddings locales y enumeran candidatos exactos para Core.
\end{description}

Core y scraper comparten una instancia MySQL, pero la propiedad se delimita por tablas y herramientas de migración: Flyway administra identidad, bundles, jobs y outbox; Alembic administra aplicaciones, fuentes y pipeline. El worker y notificaciones conservan inbox idempotentes en H2; MinIO almacena resultados temporales; PostgreSQL/pgvector pertenece al servicio semántico y puede reconstruirse desde MySQL. RabbitMQ conecta los servicios mediante mensajes que contienen identificadores y correlación, nunca URL resueltas, cookies ni secretos.

\diagramaapaisado{architecture_containers.pdf}{Contenedores implementados y dependencias principales.}{fig:contenedores-implementados}

\section{Clasificación materializada del catálogo}

El catálogo público usa tres categorías exhaustivas. Una aplicación está disponible cuando conserva una fuente \emph{direct} o \emph{fallback}, válida y con al menos un candidato \ruta{catalog\_downloadable}. Si no existe esa alternativa y alguna fuente requiere revisión manual, aparece en revisión; el resto de aplicaciones activas se clasifica como sin instalador. La antigüedad de \ruta{checked\_at} o \ruta{expires\_at} no cambia por sí sola este estado.

MySQL materializa la contribución segura de cada candidato y fuente, el \ruta{catalog\_status} de cada aplicación y una fila singleton con \ruta{total}, \ruta{available}, \ruta{review} y \ruta{missing}. Los triggers trasladan exclusivamente deltas \(+1/-1\) dentro de la transacción que modifica el dato original. La igualdad \(total=available+review+missing\) es una restricción de integridad. Una orden de mantenimiento permite auditar o reconstruir la proyección después de una restauración.

Esta proyección también actúa como versión de cambios para WebSocket. Las búsquedas y facetas filtran por el estado materializado, mientras que la selección de descarga sigue comprobando el candidato concreto, ignorando la veracidad de la proyección.

\section{Flujo general de descarga}

La figura~\ref{fig:secuencia-general} resume el recorrido desde la selección del usuario hasta la eliminación de los temporales.

\diagramaapaisado{sequenceDiagram_general.pdf}{Secuencia general de generación y descarga de un ZIP.}{fig:secuencia-general}

El backend no devuelve el ZIP en la misma petición que crea el trabajo. Tras validar la selección y las cuotas, persiste el estado inicial, sus elementos y un evento outbox. La respuesta inmediata contiene un identificador que el frontend utiliza para consultar el progreso por SSE, con sondeo como alternativa.

El worker procesa cada aplicación de forma aislada. Una incidencia no oculta el estado de las demás. Cuando hay al menos un archivo válido, el trabajo produce un ZIP parcial acompañado de errores legibles. Si no se incluye ningún archivo ni alternativa manual, el estado es \emph{FAILED}. La cancelación es idempotente y la limpieza elimina staging al terminar o fallar una fase. Core reserva almacenamiento antes de despachar el trabajo; cuando no cabe, lo mantiene en una cola FIFO sin bloquear la petición HTTP.

\imagen{download_job_states.pdf}{Máquina de estados de un trabajo de descarga.}{0.80}

\section{Resolución de instaladores}

Cada fuente combina datos declarativos y una estrategia:

\begin{itemize}
	\item Sistema operativo, arquitectura, región y país.
	\item Página oficial de partida (Winstall).
	\item Página oficial, sistema, arquitectura y evidencia de validación.
	\item Tipo de resolver y configuración.
	\item Extensiones y tipos de contenido esperados.
	\item Tiempo máximo, reintentos y número de redirecciones.
	\item Fecha y resultado de la última comprobación.
\end{itemize}

El registro de \emph{ResolverStrategy} consulta primero información estructurada de Winstall, endpoints conocidos, GitHub Releases y análisis HTML; Playwright actúa como respaldo. La arquitectura de candidatos encadena descubrimiento, transformación, clasificación de formato y validación.

El resultado del navegador no es suficiente. Un validador independiente repite las comprobaciones de HTTPS y red pública, realiza una petición controlada y verifica tamaño, MIME y hash cuando existe. Cada resolución de instalador es válida durante 24 horas desde su última comprobación correcta. Al caducar, la fuente continúa visible y seleccionable, pues la fecha por sí sola no cambia su estado público. Al preparar el ZIP, el worker solicita su resolución interna y el scraper la revalida antes de revelar la URL; si la comprobación es correcta, renueva la vigencia otras 24 horas. Un fallo terminal invalida solo ese candidato y los triggers reclasifican la aplicación; un fallo transitorio conserva el estado y permite reintentar.

\section{Generación y contenido del ZIP}

El trabajo dispone de un espacio temporal aislado. Los nombres se normalizan, se eliminan componentes de ruta y se resuelven duplicados. El contenido real es:

\begin{verbatim}
instalador-aplicacion.exe
instalador-aplicacion2.exe
manifest.json
\end{verbatim}

El manifiesto indica la fecha de creación y el resultado de cada elemento. No incluye credenciales ni URL resueltas; las incidencias se representan mediante códigos estables. El nombre de la aplicación y los metadatos necesarios permiten relacionar cada instalador con la selección.

Core no expone MinIO directamente al navegador: la ruta autorizada transmite el ZIP mediante streaming y admite \ruta{HEAD} y un rango HTTP para reanudar una descarga. El coordinador registra actividad y recibos de guardado, mantiene reservas de capacidad y elimina el resultado al confirmarse su guardado, cancelarse o agotarse los controles de inactividad. Una tarea periódica limpia estados residuales; el ciclo de vida del objeto no se basa en una expiración fija de veinticuatro horas. La lógica de entrega y reintento queda cubierta por pruebas focalizadas.

\section{Bundles}

La figura~\ref{fig:secuencia-bundles} muestra la consulta pública, el mantenimiento administrativo y la descarga de un bundle.

\diagramaapaisado{sequenceDiagram_bundle.pdf}{Secuencia de consulta, mantenimiento administrativo y descarga de un bundle.}{fig:secuencia-bundles}

La API pública lista bundles activos, consulta su detalle y crea un trabajo de descarga a partir de su composición completa. La administración crea, actualiza y borra bundles oficiales; una cuenta autenticada crea, actualiza y borra sus propios bundles públicos o privados desde el panel personal, con comprobación de propietario y versión esperada. La clonación, personalización de bundles ajenos y las estrellas no forman parte de la entrega.

Los bundles descargables no superan cien aplicaciones.

El contador de aplicaciones del bundle se mantiene junto a la entidad y se repara después de cambios autorizados. La descarga vuelve a consultar las aplicaciones, aplica el máximo y omite de forma explícita cualquier elemento no seleccionable.

\section{Instaladores Linux y alternativas manuales}

La descarga Linux admite perfiles para APT, DNF, Pacman y Zypper, además de una ruta portable. El usuario selecciona el gestor y la arquitectura; el sistema prioriza los formatos compatibles y conserva la versión solicitada cuando la fuente lo permite. Los bundles Linux incluyen instaladores automáticos, accesos \ruta{.url} a páginas oficiales HTTPS o una omisión explicada. El worker comprueba el hash SHA-256 y la firma GPG cuando existen y permite volver a una versión anterior del bundle publicado. La ausencia de instalador automático no se oculta: el trabajo termina como \emph{MANUAL\_ONLY} si todos sus elementos son alternativas manuales.

\section{Enriquecimiento de descripciones}

La figura~\ref{fig:secuencia-descripciones} refleja el flujo de generación de una descripción larga fuera de la ruta crítica.

\diagramaapaisado{sequenceDiagram_llmSearch.pdf}{Secuencia de enriquecimiento de descripciones mediante LLM.}{fig:secuencia-descripciones}

El supervisor selecciona aplicaciones sin descripción larga, calcula un hash de la entrada y encola el trabajo de forma idempotente. El descriptor reúne únicamente metadatos persistidos y llama al proveedor configurado con límites de concurrencia y tiempo.

La respuesta se valida antes de persistirla junto con idioma, proveedor, modelo, fecha y hash. El corte final permite completar por separado la descripción en español y en inglés; la falta de una de ellas no se presenta como enriquecimiento bilingüe completo. La procedencia se guarda en \ruta{long\_description\_source}: se asigna \ruta{admin\_manual} cuando el texto revisado difiere de la propuesta de IA; si coincide, se conservan el proveedor y el modelo. Esta marca clasifica el contenido, no registra las acciones de edición, por lo que un texto idéntico a la propuesta mantiene la atribución al proveedor. La respuesta normal del catálogo no expone la procedencia, aunque la previsualización administrativa muestra el origen de las sugerencias antes de publicar. El enriquecimiento automático no sobrescribe una descripción manual. Los fallos quedan registrados y pueden reintentarse. Groq y DeepSeek son intercambiables mediante configuración.

Este flujo no genera embeddings, no consulta pgvector y no interviene en la clasificación de fuentes. El indexador semántico reutiliza después la descripción persistida junto con el resto de metadatos autorizados.

\section{Búsqueda semántica y degradación segura}

El scraper publica un feed interno autenticado ordenado por identificador. Para cada aplicación activa construye el mismo texto canónico y el mismo \ruta{contentHash}; no incluye URL completas, secretos ni metadatos crudos. El indexador realiza barridos completos, actualiza únicamente documentos modificados, inactiva huérfanos y procesa embeddings mediante una cola PostgreSQL con leases.

El endpoint semántico carga solo la versión activa y enumera mediante similitud coseno. Core limita la colección y comprueba cobertura, modelo e índice: si falta alguno, si la consulta no es elegible o si aparece un error de autenticación, red o tiempo, crea un contexto literal para toda la petición. Por tanto, lista, conteo o facetas nunca mezclan un conjunto semántico parcial.

Cuando el alcance es válido, Core entrega la colección de candidatos a MySQL mediante \ruta{JSON\_TABLE}. La consulta incorpora estado, sistemas, etiquetas, editor, desempates, límite y desplazamiento sobre la tabla autoritativa. Los IDs semánticos huérfanos no producen filas porque la unión conserva la aplicación MySQL como requisito.

La interfaz presenta \emph{Literal} y \emph{Semántica}. El modo de la URL tiene prioridad sobre \ruta{localStorage}; por defecto, se elige el modo de búsqueda semántico. El fallback no modifica la preferencia guardada y muestra un aviso breve. Las respuestas públicas exponen modo solicitado y aplicado, versiones de modelo e índice y una causa opcional de degradación.

\section{Seguridad}

\subsection{Prevención de SSRF}

Las URL externas se consideran datos no confiables. Antes de cualquier conexión se validan esquema HTTPS, host y puerto; se resuelven todas las direcciones; y se rechazan redes privadas, loopback, enlace local, direcciones reservadas y puntos de metadatos de proveedores cloud. La validación se repite tras cada redirección y sobre la IP efectiva para reducir ataques de cambio de DNS.

\subsection{Validación pública de destinos}

No se confía en una allowlist de dominios como único control. La fuente aporta contexto de origen, pero la decisión final exige HTTPS, red pública, límites de redirección, tiempos de espera, MIME esperado y validación de tamaño y firma. De este modo una URL marcada como atestiguada no se considera descargable hasta comprobar el binario real.

\subsection{Límites y protección frente a abuso}

El perfil de producción versionado admite hasta 100 aplicaciones por trabajo, archivos de hasta 4 GB, cinco redirecciones y, ante fallos temporales, hasta dos reintentos de descarga (tres intentos en total). Cada trabajo descarga como máximo dos archivos en paralelo y el worker puede empaquetar dos trabajos simultáneamente. Para usuarios anónimos se configuran diez creaciones por hora y navegador y treinta por hora e IP. No se fija un máximo de 20 GiB por trabajo ni un tope de trabajos no terminales por navegador: Core reserva un máximo global de 10 GiB para temporales, ZIP y transferencias multipart, y los trabajos que no caben esperan en FIFO. El worker mantiene 8 GB libres en el host y el bucket de ZIP tiene una cuota de 10 GB. Estos valores describen la configuración versionada; no acreditan por sí solos las variables efectivas del contenedor en producción.

Los usuarios anónimos conservan la funcionalidad básica mediante una cookie de propiedad y cuotas inferiores.

\subsection{Archivos y ZIP}

La plataforma no ejecutará los instaladores. Los nombres se sanitizarán para impedir rutas absolutas, secuencias \ruta{../}, nombres reservados o colisiones. Se validará el tamaño real y se interrumpirá la descarga si supera el límite, aunque la fuente hubiera anunciado un valor menor.

\subsection{Autenticación y autorización}

Spring Security protege las operaciones por rol y propietario. El usuario solicita un enlace de un solo uso por correo; el token tiene una vigencia limitada y se consume al autenticar o crear la cuenta. El administrador conserva un acceso con contraseña almacenada como hash. Las sesiones se guardan con Spring Session JDBC y aplican cookies \emph{Secure}, \emph{HttpOnly} y \emph{SameSite}, además de CSRF. Los trabajos anónimos se protegen con una cookie aleatoria cuyo HMAC se persiste.

\section{Gestión de errores}

Los errores se clasificarán por fase:

\begin{description}
	\item[Resolución:] no existe enlace, cambió la página, se requiere interacción no soportada o el dominio final no es válido.
	\item[Validación:] protocolo, tamaño, extensión, tipo o redirección no permitidos.
	\item[Descarga:] timeout, conexión interrumpida, límite externo o archivo incompleto.
	\item[Compresión:] falta de espacio, cancelación o error al construir el ZIP.
	\item[Enriquecimiento:] proveedor no disponible, cuota agotada, respuesta inválida o contexto insuficiente.
\end{description}

El usuario recibirá una explicación comprensible y un identificador de incidencia, mientras que el registro técnico conservará detalles para el administrador en caso de que se requiera una revisión.

\section{Observabilidad}

Spring Actuator, los eventos de progreso y las métricas persistidas del scraper permiten observar:

\begin{itemize}
	\item Trabajos creados, completados, parciales, fallidos y expirados.
	\item Duración de resolución, descarga y compresión.
	\item Longitud de las colas y número de consumidores.
	\item Aplicaciones con más fallos y dominios que cambian.
	\item Espacio temporal utilizado y objetos pendientes de limpieza.
	\item Enriquecimientos de descripción, fallos del proveedor y validaciones rechazadas.
\end{itemize}

Los servicios Java exponen salud mediante Actuator y métricas compatibles, y el despliegue publica healthchecks y monitorización diaria. El scraper guarda instantáneas de colas, trabajadores y evolución; el panel recibe cambios por WebSocket y los trabajos por SSE. Estas señales permiten observar reintentos, leases y límites. No se afirma que exista un servidor Prometheus central ni una capacidad de producción medida.

\section{Riesgos y estrategia de pruebas}

El anexo de planificación evalúa riesgos inherentes y residuales en los dos cortes mediante probabilidad e impacto de 1 a 5. Se priorizan los riesgos con puntuación alta y los controles que protegen datos, integridad y disponibilidad.

La selección de casos se relaciona con requisitos funcionales y no funcionales y con características de seguridad, fiabilidad, adecuación funcional y mantenibilidad del modelo de calidad ISO/IEC 25010:2023~\cite{iso25010}. ISO/IEC 25002:2024 aporta el marco de uso de modelos de calidad~\cite{iso25002}; no se afirma conformidad certificada con esas normas. Las comprobaciones disponibles abarcan:

\begin{itemize}
	\item Pruebas unitarias de URL, dominios, IPs privadas, nombres, permisos y estados.
	\item Pruebas MySQL mediante Testcontainers para columnas generadas, triggers, rollback, concurrencia y reconciliación~\cite{testcontainers}.
	\item Páginas controladas para probar resolvers, redirecciones y descargas.
	\item Pruebas de ZIP correcto, parcial, cancelado y expirado.
	\item Pruebas de seguridad para SSRF, cuotas, autorización y datos maliciosos.
	\item Pruebas de extremo a extremo del catálogo, selección, progreso, bundles y administración.
	\item Pruebas PostgreSQL/pgvector de migraciones, dimensiones, sincronización, leases, cobertura, promoción y rollback.
	\item Pruebas de filtros, totales, facetas, páginas y fallback íntegro del modo semántico.
	\item Verificación de CPython, imports sin reactivación del GIL, Chromium y concurrencia repetida.
	\item Compilación de contratos, imágenes y documentos reproducibles.
\end{itemize}

Las pruebas locales verifican comportamientos acotados; el acceso público se observa sin lanzar ataques ni carga contra producción. La restauración completa de una copia, la purga uniforme de ciertos registros y la campaña de carga sostenida no quedan acreditadas por estas ejecuciones.

\capitulo{6}{Trabajos relacionados}

\section{Ninite}

Ninite es la referencia más cercana a la idea central del proyecto. Su página permite marcar aplicaciones y obtener un instalador personalizado que automatiza el proceso. El servicio declara que descarga las aplicaciones desde el sitio de cada editor, selecciona versiones adecuadas, omite programas ya actualizados y verifica firmas digitales o resúmenes antes de ejecutar~\cite{ninite}.

Su propuesta destaca por la sencillez: el usuario elige programas, descarga un ejecutable y lo inicia. También dispone de una oferta profesional orientada a la administración y actualización de equipos.

\section{Patch My PC Home Updater}

Patch My PC Home Updater es una aplicación gratuita para Windows orientada a instalar, actualizar y desinstalar software de uso habitual. La herramienta mantiene un catálogo de más de 500 aplicaciones y permite ejecutar instalaciones y actualizaciones silenciosas, programar comprobaciones automáticas y excluir aquellos programas que el usuario no desee actualizar. También ofrece una versión portable y una función de desinstalación múltiple~\cite{patchmypc-home-updater}.

Su propuesta está especialmente centrada en el mantenimiento continuo del equipo. Las actualizaciones se prueban para comprobar su detección e instalación y, según el proveedor, los archivos se verifican mediante VirusTotal. Además, la aplicación evita incluir software adicional no solicitado y puede impedir reinicios durante el proceso de actualización~\cite{patchmypc-home-updater}.

La principal diferencia respecto a Batch Downloader es la necesidad de ejecutar un agente local. Patch My PC administra directamente el software instalado en el equipo y automatiza operaciones que pueden requerir el cierre de aplicaciones o permisos de instalación. Batch Downloader, por el contrario, se limita a localizar y agrupar instaladores oficiales en un archivo ZIP, sin instalar, actualizar ni desinstalar programas.

\section{UBULabs}

UBULabs es un servicio de la Universidad de Burgos que proporciona a la comunidad universitaria acceso a las aplicaciones necesarias para prácticas docentes y laboratorios. El usuario accede al portal mediante sus credenciales institucionales y el catálogo disponible depende de su cuenta, del sistema operativo y de la red desde la que se conecta~\cite{ubulabs}. Algunas aplicaciones solo están disponibles desde la red de la Universidad debido a las condiciones de sus licencias.

La plataforma utiliza dos mecanismos de entrega. Mediante el \emph{streaming} de aplicaciones se despliegan programas preconfigurados desde los servidores de la Universidad para que se ejecuten en el equipo del usuario. Cuando una licencia está restringida a la red universitaria o el acceso se realiza desde el exterior, se recurre a la virtualización: la aplicación se ejecuta en un servidor y el usuario interactúa con ella mediante un cliente remoto~\cite{ubulabs}.

UBULabs comparte con Batch Downloader la idea de ofrecer un punto de acceso centralizado a un catálogo de software, pero responde a un contexto institucional. Su catálogo y sus permisos están condicionados por acuerdos de licencia, asignaturas y pertenencia a la Universidad. Batch Downloader plantea, en cambio, un catálogo público de instaladores obtenidos desde los sitios oficiales, bundles compartibles y descargas temporales, sin proporcionar licencias ni ejecutar las aplicaciones en infraestructura propia.

\section{Comparación funcional}

\begin{table}[H]
\centering
\scriptsize
\begin{tabularx}{\textwidth}{p{0.20\textwidth}XXX}
\toprule
\textbf{Característica} & \textbf{Ninite / Patch My PC} & \textbf{UBULabs} & \textbf{Batch Downloader}\\
\midrule
Finalidad principal & Instalación y mantenimiento automatizados & Acceso académico a software con licencia & Descarga conjunta de instaladores oficiales\\
Selección múltiple & Sí & No & Sí\\
Forma de entrega & Instalador o agente local & \emph{Streaming}, descarga o escritorio virtual & ZIP temporal con instaladores y manifiesto\\
Ejecución automática & Sí & Sí, local o remota según la aplicación & No; el usuario decide qué ejecutar\\
Control de acceso & Servicio general, con oferta profesional & Credenciales, permisos y condiciones de red & Catálogo público y funciones asociadas a cuentas\\
Persistencia de binarios & Gestionada por cada servicio & Aplicaciones alojadas o desplegadas por la Universidad & Prohibida por diseño para instaladores\\
Colecciones compartibles & Selección de aplicaciones & No es su objetivo & Bundles oficiales y públicos consultables y descargables\\
Participación comunitaria & Sugerencia de aplicaciones & Catálogo institucional mantenido & Bundles propios públicos o privados y mantenimiento de bundles oficiales\\
Sistemas operativos & Principalmente Windows & Depende de la aplicación y del método de acceso & Windows, Linux y macOS cuando exista instalador validado\\
Búsqueda semántica & No & No & Sí; modo semántico local con fallback literal\\
Uso de LLM & No es central & No & Solo enriquecimiento de descripciones; el ranking usa embeddings locales\\
\bottomrule
\end{tabularx}
\caption{Comparación conceptual de las soluciones relacionadas con Batch Downloader.}
\label{tabla:comparacion-trabajos-relacionados}
\end{table}

\section{Diferencias de enfoque}

Ninite y Patch My PC buscan que la instalación o actualización sea automática y silenciosa. Batch Downloader reúne archivos y conserva trazabilidad, sin ejecutar software en el equipo del usuario. Esta decisión reduce el alcance y evita desarrollar un agente local con privilegios de instalación, aunque traslada al usuario la ejecución posterior.

UBULabs resuelve un problema diferente: garantizar que una comunidad concreta pueda utilizar software sujeto a licencias institucionales. Su combinación de despliegue y virtualización permite ejecutar aplicaciones incluso cuando no pueden instalarse libremente en el equipo del usuario. Batch Downloader no cubre este escenario: incorpora software cuyos instaladores pueden obtenerse desde fuentes oficiales y no gestiona licencias de uso.

Los bundles introducen una capa organizativa. Las rutas públicas permiten listar, consultar y descargar colecciones oficiales o públicas; la administración mantiene los oficiales y cada cuenta mantiene sus propios bundles públicos o privados bajo control de propietario. La entrega no ofrece clonación, estrellas ni personalización comunitaria, aunque el esquema conserve campos heredados que podrían apoyar una ampliación futura.

La búsqueda semántica forma parte del recorrido funcional sin desplazar a MySQL. PostgreSQL, pgvector y \ruta{semantic-service} amplían la recuperación por intención, mientras que Core conserva filtros, facetas, clasificación y paginación deterministas. Groq y DeepSeek solo redactan descripciones largas; el modelo de embeddings se aprovisiona localmente, no genera recomendaciones ni crea bundles y se degrada a búsqueda literal ante cualquier incoherencia.

\section{Lecciones aplicadas al diseño}

La simplicidad de Ninite y Patch My PC demuestra la importancia de reducir pasos y explicar claramente qué ocurrirá. Batch Downloader evita que la complejidad interna de colas, workers y resolvers se traslade a la interfaz. El flujo principal es seleccionar, revisar y generar.

También resulta relevante la confianza. Ninite destaca el origen oficial y la verificación de los archivos~\cite{ninite}, mientras que Patch My PC añade pruebas previas de las actualizaciones~\cite{patchmypc-home-updater}. En Batch Downloader estas garantías se materializan mediante validación HTTPS y de red pública tras las redirecciones, manifiesto, fechas de comprobación y mensajes explícitos.

UBULabs muestra, además, que la disponibilidad de una aplicación puede depender del usuario, la plataforma, la ubicación de red y la licencia. Aunque Batch Downloader no gestiona licencias institucionales, su catálogo representa los sistemas compatibles y distingue de forma exhaustiva \emph{Disponible}, \emph{Revisión} y \emph{Sin instalador}. La antigüedad de una resolución no cambia por sí sola esa categoría: la revalidación se realiza justo antes de descargar.

Por último, la comparación ayuda a delimitar el proyecto. Batch Downloader no pretende sustituir un gestor de paquetes del sistema operativo ni automatizar instalaciones con privilegios. Su aportación se sitúa en la preparación de una descarga conjunta, los bundles administrados y un catálogo controlado con clasificación transaccional y recuperación semántica degradable.

\capitulo{7}{Conclusiones y líneas de trabajo futuras}

\section{Conclusiones de la implementación}

La implementación documentada confirma que Batch Downloader es más que una página con enlaces. La generación conjunta de instaladores exige resolver fuentes cambiantes, validar destinos, procesar archivos grandes de forma asíncrona, comunicar progreso y eliminar temporales. Estos requisitos justifican la separación entre API, workers, resolvers y almacenamiento temporal.

La decisión de no conservar instaladores permanentemente se mantiene como principio central. El catálogo almacena conocimiento sobre las aplicaciones y sus fuentes, pero cada archivo se obtiene en el contexto de un trabajo. Esta política reduce duplicación y favorece la actualización, aunque obliga a asumir que las páginas externas pueden cambiar y que es necesario mantener resolvers y registros.

Los bundles aportan reutilización sin perder control sobre su composición. Las rutas públicas permiten listarlos, consultar su detalle y crear un trabajo; las rutas administrativas concentran su mantenimiento. La clonación, las estrellas y la personalización comunitaria no están implementadas y no se presentan como resultado de esta entrega.

La búsqueda semántica se incorpora sin convertir el modelo en autoridad del catálogo. El feed canónico, el indexador, pgvector, la búsqueda exacta y la fusión RRF forman una ruta funcional con fallback literal íntegro. Los pesos se ejecutan localmente; Groq y DeepSeek se emplean exclusivamente para redactar descripciones largas y no intervienen en el ranking, no aprueban enlaces, no recomiendan aplicaciones y no modifican configuraciones críticas.

La arquitectura se ha incorporado por etapas. La prioridad fue validar un núcleo seguro y observable antes de introducir el LLM en el enriquecimiento de texto. La proyección incremental del catálogo completa ese núcleo: los triggers de MySQL propagan cambios de candidato, fuente y aplicación hasta una fila singleton, manteniendo \emph{Disponible}, \emph{Revisión} y \emph{Sin instalador} como partición exhaustiva de las aplicaciones activas.

\section{Limitaciones actuales}

La tabla~\ref{tabla:resultados-verificacion-semantica} recoge la ejecución local del 23 de julio de 2026 sobre el árbol de trabajo. Estas cifras describen pruebas funcionales, no una campaña de carga ni una evaluación de calidad de los modelos sobre un conjunto humano definitivo.

\begin{table}[H]
\centering
\small
\begin{tabularx}{\textwidth}{p{0.28\textwidth}p{0.20\textwidth}X}
\toprule
\textbf{Ámbito} & \textbf{Resultado} & \textbf{Cobertura relevante}\\
\midrule
Core API & 41 pruebas superadas & Contrato híbrido, fallback, RRF, filtros y consultas de catálogo\\
Frontend & 51 pruebas y build & Selector, precedencia URL/\ruta{localStorage}, aviso y tipos\\
Scraper & 156 pruebas superadas & Feed canónico, hash, MySQL~8.4, proyección y driver aiomysql\\
Servicio semántico & 10 pruebas superadas & Migraciones pgvector, dimensiones, sincronización, promoción, variantes de entrenamiento y HNSW exacto\\
Compose & Stack construido y operativo & Healthchecks, migraciones, 13.667 embeddings E5 y fallback íntegro a través de Nginx\\
Entrenamiento acotado & 3 modelos y 25 variantes & Snapshot de 512 aplicaciones, artefactos LoRA recargables e informes JSON, CSV y Markdown\\
CPython 3.14/3.14t & 5 repeticiones por configuración & Checksums equivalentes; aceleración a 8 hilos de \(3{,}100\times\) en HTML y \(3{,}380\times\) en HTTP controlado\\
\bottomrule
\end{tabularx}
\caption{Resultados funcionales obtenidos durante la verificación integral local del stack.}
\label{tabla:resultados-verificacion-semantica}
\end{table}

El perfil \ruta{--smoke} creó el snapshot \ruta{06e6b8dda8bc}, con 512 aplicaciones y 2.830 consultas generadas, y acotó la evaluación efectiva a 256 documentos y ocho consultas de validación. En ese alcance de integración, BGE-M3 \emph{zero-shot} obtuvo nDCG@10 y MRR@10 iguales a \(1{,}0000\), mientras el literal obtuvo \(0{,}9126\) y \(0{,}8889\), respectivamente. Estas cifras no constituyen una comparación de calidad concluyente: el conjunto es pequeño y sintético, y todas las variantes se marcaron expresamente como no elegibles. Por ello no se consultó test ni se promocionó un modelo; E5-base \emph{zero-shot} continuó activo.

La comparación del runtime utilizó 200 tareas por configuración, HTML almacenado, un servidor HTTP controlado, 1/2/4/8 hilos y cinco repeticiones. A un hilo, 3.14t asumió una sobrecarga en la carga puramente CPU (\(0{,}916\times\)); a ocho hilos alcanzó \(699{,}15\) tareas/s frente a \(225{,}56\) en el control estándar. En la carga HTTP controlada obtuvo \(642{,}49\) frente a \(190{,}11\) tareas/s. Los checksums fueron idénticos en todas las configuraciones.

La selección automática de producción sigue exigiendo ejecutar el perfil completo: el ajuste debe superar a su variante \emph{zero-shot} y no reducir MRR@1 respecto al literal. Solo entonces se confirma una vez sobre test y se activa con cobertura completa; si no existe un ganador elegible, se conserva E5-base \emph{zero-shot}.

La compatibilidad con Windows, Linux y macOS depende de la disponibilidad y condiciones de las fuentes oficiales. La interfaz muestra únicamente sistemas con una fuente seleccionable. \ruta{checked\_at} y \ruta{expires\_at} no retiran por sí solos un candidato estructuralmente válido: si ha caducado, se revalida bajo bloqueo justo antes de descargar. Un fallo terminal invalida el candidato y activa la reclasificación por triggers; uno transitorio conserva la categoría y permite reintentar.

La proyección incluye un comando de auditoría y otro de reparación para restauraciones manuales. Su corrección transaccional debe seguir comprobándose con MySQL 8.4 real, escenarios multifuente, movimientos, borrados, rollback y concurrencia; una prueba basada solo en fragmentos SQL no cubre esas propiedades.

Los servicios se ejecutan con JDK 26, pero el proceso Maven fija la generación de bytecode Java 25 para mantener compatibilidad de herramientas. La matriz de dependencias sigue siendo una responsabilidad de mantenimiento: cada actualización de Spring Boot, SonarJava o la plataforma Java debe validarse en integración continua antes de su adopción.

\section{Líneas de trabajo futuras}

\subsection{Rendimiento y pruebas de extremo a extremo}

La siguiente línea de trabajo es ejecutar una campaña reproducible con MySQL, RabbitMQ, MinIO, Mailpit y un servidor HTTP controlado. Debe cubrir el flujo visitante--bundle--job--ZIP, revalidación JIT, cancelación, resultados parciales, recuperación tras caída y los límites de 100 aplicaciones y 20 GiB. Las pruebas de SSRF y limpieza temporal seguirán siendo criterios de aceptación, no mejoras opcionales.

\subsection{Resolvers y salud del catálogo}

Se podrán desarrollar resolvers específicos para proveedores frecuentes y mejorar el panel que prioriza aplicaciones rotas. El historial permitirá detectar cambios repetidos y decidir qué entradas necesitan mantenimiento manual. El estado público continuará dependiendo de invalidaciones explícitas, no del paso del tiempo.

\subsection{Rendimiento y streaming}

Tras medir el uso de disco y memoria se estudiará añadir archivos al ZIP mediante streaming. Esta técnica reduciría almacenamiento temporal, pero requiere resolver cuidadosamente los fallos parciales, la cancelación y la imposibilidad de conocer el tamaño final.

\subsection{Evaluación semántica continua}

La ruta funcional queda preparada para repetir snapshots y benchmarks sin cambiar la API. Como continuación se deberá ampliar el conjunto de consultas humanas en español, errores tipográficos, sinónimos e intención cruzada entre idiomas; vigilar deriva del catálogo; y revisar periódicamente umbral, peso RRF y recall HNSW. La promoción seguirá requiriendo cobertura completa y conservará la versión anterior para rollback.

\subsection{Modelos locales, eficiencia y privacidad}

Los tres modelos de embeddings ya se ejecutan localmente. Como trabajo posterior se estudiarán cuantización, aceleradores, batching adaptativo y modelos más pequeños, comparando calidad, latencia y consumo energético. Las consultas y descripciones continuarán fuera de servicios de inferencia externos y los informes no incluirán secretos ni datos de contacto.

\subsection{Despliegue y escalado}

Docker Compose cubre el entorno inicial con MySQL compartido, H2 local en los consumidores, RabbitMQ, MinIO y Mailpit. Un despliegue distribuido deberá añadir gestión de secretos, límites de recursos, copias de seguridad, políticas de red y una solución de métricas; Prometheus y Kubernetes no forman parte de esta entrega.

\subsection{Accesibilidad e internacionalización}

La interfaz podrá ampliarse con varios idiomas, navegación completa por teclado, mensajes compatibles con lectores de pantalla y pruebas sobre criterios WCAG. La accesibilidad deberá validarse desde las primeras pantallas y no añadirse únicamente al final.

\section{Cierre}

Batch Downloader reúne en una experiencia sencilla un conjunto complejo de responsabilidades: catálogo, seguridad, procesamiento asíncrono, bundles administrados, enriquecimiento de descripciones y búsqueda híbrida. La implementación establece límites claros para cada componente y proporciona una base verificable para continuar el trabajo sin presentar como terminadas las funciones comunitarias, una campaña de carga de producción o el despliegue público.



\bibliographystyle{plain}
\bibliography{bibliografia}

\end{document}
